\documentclass{iopconfser}

\usepackage{graphicx}
\usepackage{amsmath}
\usepackage{booktabs}

\begin{document}

\title{Grounded Vision--Language Adaptation on Chest CT and Reports for
Clinically-Oriented Zero-Shot Pulmonary Nodule Detection and Case Retrieval}

\author{Dionisius Seraf Saputra$^{1}$, Nafis Naufal Rahman$^{2}$}

\affil{$^1$Department of Information System, Universitas Brawijaya, Malang, Indonesia}
\affil{$^2$Department of Informatics Engineering, Universitas Brawijaya, Malang, Indonesia}

\email{dionisiusve32@student.ub.ac.id, nafisnaufal1426@student.ub.ac.id}

\begin{abstract}
Manual annotation of 3D medical modalities such as Chest Computed Tomography
(CT) is highly labour-intensive, limiting the scalability of deep learning
systems for thoracic diagnosis. While Vision-Language Models (VLMs) show
significant promise in 2D imaging, applying them natively to volumetric CT
data remains challenging due to memory constraints and the difficulty of
fine-grained spatial localisation. This paper presents a parameter-efficient
adaptation framework built on VILA-M3, an expert-routed medical VLM developed
under the MONAI project. The framework is adapted to thoracic clinical tasks
via Low-Rank Adaptation (LoRA) on the large-scale CT-RATE dataset, aligning
the model's language representations with chest CT visual patterns and radiology
report semantics. A built-in dynamic routing mechanism delegates high-resolution
volumetric localisation to VISTA3D, a specialist 3D segmentation expert, while
the VLM retains semantic control. Zero-shot transferability of the adapted model
is evaluated on the LIDC-IDRI reference database for pulmonary nodule
localisation and on a CT-RATE held-out split for natural language-based case
retrieval. This work demonstrates the practical viability of expert-routed VLM
adaptation as a scalable approach to clinical decision support that substantially
reduces reliance on densely annotated training data.
\end{abstract}


% -----------------------------------------------------------------------
\section{Introduction}
% -----------------------------------------------------------------------

Chest CT is a primary imaging modality for diagnosing thoracic diseases
including pulmonary nodules, pneumonia, pleural effusion, and cardiovascular
conditions. Despite advances in deep learning, high-performing models in this
domain still depend on large datasets annotated with dense bounding boxes or
pixel-level segmentation masks---an expensive process requiring considerable
radiological expertise and time~\cite{litjens2017}.

Vision-Language Models (VLMs) offer an alternative by learning visual
representations from the natural pairing of medical images and clinical
free-text reports, bypassing the need for explicit spatial labels~\cite{clip2021,
medclip2022}. Early medical VLMs such as MedCLIP~\cite{medclip2022} and
BioViL-T~\cite{biovilt2023} have demonstrated strong performance on chest X-ray
classification and report grounding. However, extending these approaches to 3D CT
volumes introduces fundamental challenges: processing full volumetric sequences of
tokens is computationally prohibitive, and most VLM architectures lack mechanisms
for precise 3D spatial reasoning~\cite{vilaM32024}.

VILA-M3~\cite{vilaM32024} addresses this problem by introducing an expert routing
architecture within a VLM pipeline. Rather than processing raw 3D voxels, the VLM
acts as an orchestrator that delegates spatial tasks to specialist expert models via
the MONAI VLM Agent Framework~\cite{monai2022}. This enables high-fidelity 3D
analysis without native volumetric tokenisation. Despite this architectural
advantage, VILA-M3 is pre-trained on general medical image-text data and has not
been specifically adapted for thoracic CT interpretation.

This paper proposes a targeted instruction-tuning pipeline that fine-tunes VILA-M3
on CT-RATE~\cite{ctrate2024}, a large-scale dataset of paired chest CT volumes and
radiology reports. Adaptation is achieved efficiently via LoRA~\cite{lora2022},
requiring only a fraction of the parameters of full fine-tuning. The framework's
routing mechanism directs spatial localisation queries to VISTA3D~\cite{vista3d2024},
a versatile 3D CT segmentation expert. Zero-shot transfer is then evaluated on
LIDC-IDRI~\cite{lidc2011}, an entirely unseen benchmark of expert-annotated
pulmonary nodules, alongside a case retrieval evaluation on the CT-RATE hold-out
split.

\subsection*{Key Contributions}

\begin{enumerate}
  \item \textbf{Thoracic Instruction-Tuning Pipeline.} A parameter-efficient
    fine-tuning strategy using LoRA that aligns VILA-M3's representations with
    thoracic clinical language and chest CT visual patterns derived from CT-RATE
    report pairs, enabling both detection and retrieval task conditioning.

  \item \textbf{Expert Routing Evaluation.} Systematic assessment of the MONAI
    VLM Agent Framework's dynamic routing mechanism, which delegates volumetric
    spatial tasks to VISTA3D while preserving natural language grounding in the
    VLM backbone.

  \item \textbf{Zero-Shot Cross-Dataset Validation.} A rigorous evaluation
    protocol on LIDC-IDRI for pulmonary nodule localisation and on a CT-RATE
    held-out split for text-to-image case retrieval, benchmarked against
    supervised detection models and retrieval baselines.
\end{enumerate}


% -----------------------------------------------------------------------
\section{Related Work}
% -----------------------------------------------------------------------

\subsection{Vision-Language Models in Radiology}

Contrastive image-text pretraining, established by CLIP~\cite{clip2021}, has
been adapted to the medical domain by methods such as MedCLIP~\cite{medclip2022}
and BioViL-T~\cite{biovilt2023}. These approaches learn representations from
unpaired or temporally paired radiology report-image corpora and achieve
competitive zero-shot classification performance on chest X-ray benchmarks.
More recent large multimodal models such as GPT-4V~\cite{gpt4_2023} and
Med-Gemini extend report-grounded reasoning to open-ended clinical question
answering. However, these models remain constrained to 2D input modalities and
lack mechanisms for volumetric spatial grounding.

\subsection{3D Medical Image Analysis}

Volumetric CT analysis has been dominated by convolutional architectures.
nnU-Net~\cite{nnunet2021} established a robust self-configuring baseline for 3D
segmentation across diverse clinical tasks. For pulmonary nodule detection
specifically, 3D Faster R-CNN-based systems and sliding-window CNNs have set
strong supervised benchmarks on LIDC-IDRI~\cite{lidc2011}. VISTA3D~\cite{vista3d2024}
extends this paradigm with a promptable 3D segmentation model that accepts text
and point prompts, making it well-suited as a spatially-precise expert in a
VLM routing architecture.

\subsection{Expert-Routed and Agent-Based Medical VLMs}

VILA-M3~\cite{vilaM32024} builds on the VILA visual language architecture~\cite{vila2024}
and the LLaMA language backbone~\cite{llama2023} to introduce expert routing:
the VLM's language head generates structured routing signals that invoke
specialist models for tasks requiring high-resolution spatial reasoning. The
MONAI VLM Agent Framework~\cite{monai2022} provides the infrastructure for
this orchestration. Our work extends VILA-M3 to thoracic CT by fine-tuning the
routing and language components on domain-specific data, and performs the first
systematic zero-shot evaluation of the resulting system on an external nodule
detection benchmark.


% -----------------------------------------------------------------------
\section{Proposed Methodology}
% -----------------------------------------------------------------------

\subsection{Dataset Configuration}

The framework employs a dual-dataset strategy that strictly separates
adaptation from evaluation:

\begin{itemize}
  \item \textbf{CT-RATE}~\cite{ctrate2024}: A large-scale dataset of over
    50{,}000 reconstructed 3D chest CT volumes paired with structured radiology
    reports. CT-RATE serves as the training corpus for instruction-tuning. A
    randomly sampled 10\% held-out split is reserved exclusively for case
    retrieval evaluation.

  \item \textbf{LIDC-IDRI}~\cite{lidc2011}: A reference database of 1{,}018
    CT scans with expert-consensus bounding box annotations for pulmonary
    nodules, contributed by four independent radiologists. This dataset is
    withheld entirely from training and used only for zero-shot localisation
    benchmarking.
\end{itemize}

\subsection{Model Architecture and Dynamic Routing}

The framework positions VILA-M3 as an intelligent orchestrator between the
clinical query and specialist processing pipelines. Figure~\ref{fig:architecture}
provides an overview of the system.

\begin{figure}[htbp]
  \centering
  \includegraphics[width=0.92\columnwidth]{figures/architecture.pdf}
  \caption{Overview of the proposed framework. VILA-M3 receives multimodal
    input (key CT slices and clinical query) and classifies query intent. Localisation
    queries are routed to VISTA3D via the MONAI Agent Framework; retrieval queries
    are resolved through embedding-space matching. The LoRA fine-tuning stage
    (CT-RATE) is applied to both pathways.}
  \label{fig:architecture}
\end{figure}

The routing pipeline proceeds in three stages:

\begin{enumerate}
  \item \textbf{Global Context Extraction.} VILA-M3 receives the user query
    (e.g., \textit{``Localise the pulmonary nodule in this CT scan''}) together
    with a compact volumetric representation of the CT input, constructed by
    sampling key axial slices at uniform anatomical intervals.

  \item \textbf{Task Classification and Routing.} The model classifies the
    query intent into one of two pathways via the MONAI VLM Agent Framework.
    Localisation queries trigger delegation to VISTA3D; retrieval queries are
    handled internally through the VLM's embedding space.

  \item \textbf{Expert Execution and Response Fusion.} For localisation queries,
    VISTA3D processes the full-resolution CT volume and returns voxel-level
    segmentation masks and 3D bounding coordinates. The VLM integrates these
    spatial outputs with textual context to produce a grounded, language-anchored
    response. For retrieval queries, the VLM embedding is compared against the
    CT-RATE corpus index using cosine similarity.
\end{enumerate}

This separation of concerns preserves the VLM's semantic reasoning capacity
while delegating computationally expensive volumetric processing to a
purpose-built expert, avoiding the memory infeasibility of native 3D token
processing.

\subsection{Instruction Fine-Tuning Strategy}

We apply Low-Rank Adaptation (LoRA)~\cite{lora2022} to the VILA-M3 8B
checkpoint (\texttt{MONAI/Llama3-VILA-M3-8B}), targeting the query and value
projection matrices in the attention layers. The LoRA rank is set to
$r = 16$ with a scaling factor $\alpha = 32$. Training data is constructed
from the CT-RATE training split by converting each CT volume and paired report
into two instruction types:

\begin{itemize}
  \item \textbf{Detection instructions.} Input: a prompt of the form
    \textit{``Identify and localise thoracic abnormalities in the provided CT
    scan.''} Output: a structured description referencing anatomical regions
    and abnormality types extracted from the paired radiology report.

  \item \textbf{Retrieval instructions.} Input: a prompt of the form
    \textit{``Retrieve the most similar case to: [report excerpt].''} Output:
    the corresponding CT embedding identifier, with hard-negative sampling
    from CT-RATE to ensure discriminative training.
\end{itemize}

Training is performed on a single NVIDIA L40 GPU (48\,GB VRAM). The VILA-M3
8B model occupies approximately 18\,GB; VISTA3D requires an additional 12\,GB,
accommodating both within a single GPU during inference. The AdamW optimiser is
used with a learning rate of $1 \times 10^{-4}$, a cosine annealing schedule,
and an effective batch size of 4. Training runs for 3 epochs over the CT-RATE
training split.


% -----------------------------------------------------------------------
\section{Experimental Setup and Evaluation}
% -----------------------------------------------------------------------

\subsection{Baseline Models}

The adapted framework is compared against two categories of baselines:

\begin{enumerate}
  \item \textbf{Detection baselines.} Supervised 3D models trained on
    LIDC-IDRI annotations directly: nnU-Net~\cite{nnunet2021} and a 3D
    Faster R-CNN variant. These represent an upper bound achievable with
    full annotation access and establish the cost of annotation-free
    generalisation.

  \item \textbf{Retrieval baselines.} Zero-shot retrieval using the
    pre-fine-tuned VILA-M3 8B checkpoint (no CT-RATE adaptation) and
    MedCLIP~\cite{medclip2022}, representing the performance obtainable
    from generalist medical representations without thoracic-specific
    instruction tuning.
\end{enumerate}

\subsection{Evaluation Metrics}

\textbf{Zero-Shot Pulmonary Nodule Localisation (LIDC-IDRI).}
Performance is measured using Intersection over Union (IoU) and mean Average
Precision (mAP) computed against the consensus bounding box annotations. These
metrics penalise spatial inaccuracies rigorously, ensuring that predictions
reflect genuine localisation rather than coarse presence detection.

\textbf{Text-to-Image Case Retrieval (CT-RATE hold-out).}
Performance is measured using Recall@$K$ for $K \in \{1, 5, 10\}$. In clinical
triage, retrieving the correct historical case within the top-$K$ results
translates directly to workflow efficiency; Recall@$K$ therefore provides a
practical utility measure for semantic case search.

\subsection{Expected Outcomes}

The fine-tuned framework is expected to demonstrate meaningful zero-shot
localisation improvements over the pre-fine-tuned VILA-M3 baseline on LIDC-IDRI,
while remaining below supervised models trained on LIDC-IDRI annotations---a
gap that quantifies the residual cost of annotation-free transfer. For case
retrieval, Recall@$K$ improvements over MedCLIP and the pre-fine-tuned baseline
on the CT-RATE hold-out split would confirm that thoracic instruction tuning
yields clinically useful semantic embeddings.


% -----------------------------------------------------------------------
\section{Conclusion}
% -----------------------------------------------------------------------

This paper presents a parameter-efficient adaptation framework for expert-routed
vision-language models applied to 3D chest CT analysis. By fine-tuning VILA-M3 8B
on CT-RATE using LoRA and leveraging the MONAI VLM Agent Framework's built-in
VISTA3D routing, the framework addresses the fundamental challenge of applying
VLMs to volumetric medical data without native 3D processing. Evaluation on the
zero-shot LIDC-IDRI benchmark and the CT-RATE retrieval hold-out assesses
whether thoracic instruction tuning yields clinically meaningful improvements
over the pre-fine-tuned baseline. This approach has the potential to assist
radiologists in initial triage and semantic case retrieval, reducing dependence
on costly dense annotations. Future work will investigate multi-expert routing
strategies and extension to additional thoracic pathologies beyond pulmonary
nodules.


% -----------------------------------------------------------------------
\section*{Acknowledgements}
% -----------------------------------------------------------------------

The authors thank the MONAI Project and NVIDIA for publicly releasing the
VILA-M3 framework and model checkpoints under open licences.


% -----------------------------------------------------------------------
\begin{thebibliography}{99}
% -----------------------------------------------------------------------

\bibitem{litjens2017}
Litjens G, Kooi T, Bejnordi B E, Setio A A A, Ciompi F, Ghafoorian M,
van der Laak J A W M, van Ginneken B and S\'{a}nchez C I 2017
A survey on deep learning in medical image analysis
\textit{Med. Image Anal.} \textbf{42} 60--88

\bibitem{clip2021}
Radford A, Kim J W, Hallacy C, Ramesh A, Goh G, Agarwal S, Sastry G,
Askell A, Mishkin P, Clark J, Krueger G and Sutskever I 2021
Learning transferable visual models from natural language supervision
\textit{Proc. 38th Int. Conf. on Machine Learning (ICML)} 8748--8763

\bibitem{medclip2022}
Wang Z, Wu Z, Agarwal D and Sun J 2022
MedCLIP: Contrastive learning from unpaired medical images and text
\textit{Proc. 2022 Conf. on Empirical Methods in Natural Language Processing
(EMNLP)} 3876--3887

\bibitem{biovilt2023}
Bannur S, Hyland S L, Liu Q, Perez-Garcia F, Ilse M, Castro D C, Boecking B,
Sharma H, Bitterman T, Schwaighofer A, Wetscherek M T, Lungren M, Nori A,
Alvarez-Valle J and Oktay O 2023
Learning to exploit temporal structure for biomedical vision-language processing
\textit{Proc. IEEE/CVF Conf. on Computer Vision and Pattern Recognition (CVPR)}
15016--15027

\bibitem{gpt4_2023}
Achiam J, Adler S, Agarwal S et al.\ 2023
GPT-4 technical report
\textit{arXiv preprint} arXiv:2303.08774

\bibitem{nnunet2021}
Isensee F, Jaeger P F, Kohl S A A, Petersen J and Maier-Hein K H 2021
nnU-Net: a self-configuring method for deep learning-based biomedical image
segmentation
\textit{Nat. Methods} \textbf{18} 203--211

\bibitem{vista3d2024}
He Y, Yang D, Nath V, Myronenko A, Xu Z, Li W, Tang Y, Zhao C, Huang R
and Xu D 2024
VISTA3D: Versatile imaging segmentation and annotation model for 3D computed
tomography
\textit{arXiv preprint} arXiv:2406.05285

\bibitem{monai2022}
Cardoso M J, Li W, Brown R, Ma N, Kerfoot E, Wang Y, Murrey B, Myronenko A,
Zhao C, Yang D, Nath V, He Y, Xu Z, Hatamizadeh A, Zhu W, Liu Y, Zheng M,
Tang Y, Yang I, Yao V, Billeh Y, Zhang D, Foo B, Feng A, Young B, Larson E
and Xu D 2022
MONAI: An open-source framework for deep learning in healthcare
\textit{arXiv preprint} arXiv:2211.02701

\bibitem{vilaM32024}
Nath V, Li W, Yang D, Myronenko A, Zheng M, Lu Y, Liu Z, Yin H, Law Y M,
Tang Y and Xu D 2024
VILA-M3: Enhancing vision-language models with medical expert knowledge
\textit{arXiv preprint} arXiv:2411.12915

\bibitem{vila2024}
Lin J, Yin H, Ping W, Lu Y, Molchanov P, Tao A, Mao H, Kautz J,
Shoeybi M and Han S 2024
VILA: On pre-training for visual language models
\textit{Proc. IEEE/CVF Conf. on Computer Vision and Pattern Recognition (CVPR)}

\bibitem{llama2023}
Touvron H, Martin L, Stone K, Albert P, Almahairi A, Babaei Y,
Bashlykov N, Batra S, Bhargava P, Bhosale S, Bikel D, Blecher L and
Scialom T 2023
Llama 2: Open foundation and fine-tuned chat models
\textit{arXiv preprint} arXiv:2307.09288

\bibitem{lora2022}
Hu E J, Shen Y, Wallis P, Allen-Zhu Z, Li Y, Wang S, Wang L and Chen W 2022
LoRA: Low-rank adaptation of large language models
\textit{Proc. 10th Int. Conf. on Learning Representations (ICLR)}

\bibitem{ctrate2024}
Hamamci I E, Er B, Alber G, Simsek O, Simsar E, Senol E, Nalci M A,
Gokberk B, Hayat R S, Yuksel T and Catli S 2024
A foundation model utilizing chest CT volumes and radiology reports for
supervised-level zero-shot detection of abnormalities
\textit{arXiv preprint} arXiv:2403.17834

\bibitem{lidc2011}
Armato S G, McLennan G, Bidaut L, McNitt-Gray M F, Meyer C R, Reeves A P,
Zhao B, Aberle D R, Henschke C I, Hoffman E A, Kazerooni E A, MacMahon H,
Van Beeke E J R, Yankelevitz D and Petrick N 2011
The Lung Image Database Consortium (LIDC) and Image Database Resource
Initiative (IDRI): a completed reference database of lung nodules on CT scans
\textit{Med. Phys.} \textbf{38} 915--931

\end{thebibliography}

\end{document}
